\documentclass[UTF8,zihao=-4,fontset=none]{ctexart}
\usepackage[a4paper,margin=2.2cm]{geometry}
\usepackage{amsmath,amssymb,bm}
\usepackage{booktabs,longtable,array,multirow}
\usepackage{tikz}
\usetikzlibrary{arrows.meta,positioning,fit,shapes.geometric}
\usepackage{xcolor}
\usepackage{listings}
\usepackage{enumitem}
\usepackage{hyperref}
\usepackage{caption}
\usepackage{float}
\setCJKmainfont{Noto Serif CJK SC}
\setCJKsansfont{Noto Sans CJK SC}
\hypersetup{colorlinks=true,linkcolor=blue,urlcolor=blue,citecolor=blue}
\setlength{\parindent}{2em}
\setlength{\parskip}{0.25em}
\renewcommand{\arraystretch}{1.25}
\lstset{
  language=Python,
  basicstyle=\ttfamily\small,
  keywordstyle=\bfseries,
  commentstyle=\itshape,
  stringstyle=\ttfamily,
  showstringspaces=false,
  breaklines=true,
  frame=single,
  columns=fullflexible,
  keepspaces=true,
  tabsize=4
}
\title{层次分析法（AHP）评价类问题建模规范\\\large 理论原理、建模步骤、结论评价与 Python 代码}
\author{适用于数学建模论文写作模板}
\date{2026年7月}

\begin{document}
\maketitle
\begin{abstract}
层次分析法（Analytic Hierarchy Process, AHP）适合处理“多个评价指标、多个备选方案、需要给出综合排名或等级”的评价类问题。本文给出一套可直接写进建模论文的规范：先明确评价目标和方案，再建立目标层、准则层、方案层的层级结构；随后构造判断矩阵并检验一致性；最后计算指标权重、方案得分、总排序及总排序一致性。文末给出 Python 代码模板，可用于计算权重、一致性比率、方案得分和排序。
\end{abstract}

\tableofcontents
\newpage

\section{方法适用场景与论文写法总流程}
AHP 的核心思想是：把复杂评价问题拆成若干层次，在同一层内部只做“两两比较”，再把局部判断合成为整体权重。因此它特别适合“评价、选择、排序、优选、绩效打分、风险评估、环境质量评价”等问题。

论文中可以按以下六步组织：
\begin{enumerate}[label=Step\arabic*：,leftmargin=3.2em]
  \item \textbf{解决评价类问题。} 明确评价对象、评价目标、备选方案和评价指标，说明 AHP 适用的原因。
  \item \textbf{画出层级结构图。} 至少包含目标层、准则层、方案层；如果指标较多，可增加子准则层。
  \item \textbf{构造判断矩阵。} 对同一层指标进行两两比较，确定评价指标孰重孰轻。
  \item \textbf{依照评价指标对各方案打分。} 可以用定量数据归一化，也可以让专家依据 1--9 标度或等级标准评分。
  \item \textbf{求出权重并填表。} 计算准则权重、方案局部权重或指标评分，合成最终得分。
  \item \textbf{层次总排序一致性检验。} 若一致性不合格，应返回重新设计判断矩阵；合格后再给出结论评价。
\end{enumerate}

\section{理论原理：为什么这种评价可以用于实际}
\subsection{基本假设}
AHP 能够用于实际评价，通常需要满足以下假设：
\begin{enumerate}[label=H\arabic*：,leftmargin=3em]
  \item \textbf{可分解性假设。} 总目标可以分解为若干准则或子准则，各准则能够共同解释评价对象的优劣。
  \item \textbf{同层可比较假设。} 同一层的因素相对于上一层目标具有可比性，可以回答“$C_i$ 相比 $C_j$ 重要多少”。
  \item \textbf{正互反假设。} 判断矩阵 $A=(a_{ij})$ 满足 $a_{ij}>0$，$a_{ii}=1$，$a_{ji}=1/a_{ij}$。
  \item \textbf{近似一致性假设。} 专家判断不要求完全一致，但不能自相矛盾过大。例如若 $C_1$ 比 $C_2$ 重要 3 倍，$C_2$ 比 $C_3$ 重要 2 倍，则 $C_1$ 对 $C_3$ 的重要性应接近 6 倍。
  \item \textbf{加权合成假设。} 当各指标已经同向化、无量纲化后，方案综合效用可近似写成加权和：
  \begin{equation}
    S_k=\sum_{i=1}^{n} w_i x_{ik}, \quad \sum_{i=1}^{n}w_i=1,\quad w_i\ge 0 .
  \end{equation}
\end{enumerate}

\subsection{判断矩阵与特征向量原理}
设评价指标真实权重为 $\bm w=(w_1,w_2,\ldots,w_n)^T$。如果专家判断完全一致，则两两比较结果应满足
\begin{equation}
  a_{ij}=\frac{w_i}{w_j},\qquad a_{ij}a_{jk}=a_{ik} .
\end{equation}
此时判断矩阵为
\begin{equation}
A=\begin{pmatrix}
1 & w_1/w_2 & \cdots & w_1/w_n\\
w_2/w_1 & 1 & \cdots & w_2/w_n\\
\vdots & \vdots & \ddots & \vdots\\
w_n/w_1 & w_n/w_2 & \cdots & 1
\end{pmatrix}.
\end{equation}
直接计算可得
\begin{equation}
  A\bm w=n\bm w .
\end{equation}
因此在完全一致时，权重向量就是判断矩阵对应最大特征值 $\lambda_{\max}=n$ 的特征向量。实际专家判断常有误差，所以取最大特征根对应的正特征向量并归一化：
\begin{equation}
  A\bm w=\lambda_{\max}\bm w,
  \qquad \sum_{i=1}^n w_i=1.
\end{equation}

\subsection{概率统计解释：专家打分为什么能推广到实际}
AHP 中的专家判断可以理解为对“真实相对重要程度”的带误差观测。常用的统计解释是：
\begin{equation}
  \log a_{ij}=\log w_i-\log w_j+\varepsilon_{ij},
\end{equation}
其中 $\varepsilon_{ij}$ 表示专家认知误差、信息误差或主观波动。若误差均值接近 0，且多位专家的判断相互独立或弱相关，则对判断矩阵取几何平均可以降低随机误差：
\begin{equation}
  \bar a_{ij}=\left(\prod_{r=1}^{q} a_{ij}^{(r)}\right)^{1/q}.
\end{equation}
在这个意义下，AHP 不是把某一个人的主观意见当成客观事实，而是把专家群体判断转化为对潜在权重的估计。若判断矩阵通过一致性检验，说明这些估计没有明显逻辑矛盾；若样本数据经过同向化和标准化，方案得分就可以被解释为在既定评价体系下的综合效用估计。

需要注意：AHP 的结论只在“指标体系合理、专家来源可靠、判断矩阵一致性通过、评分数据可比”的前提下成立。它给出的是一种可解释、可复核的决策依据，而不是绝对真理。

\section{层级结构图：目标层、准则层、方案层}
论文中通常先给出类似图 \ref{fig:hierarchy} 的结构图。若准则层下面还有二级指标，则在准则层与方案层之间增加“子准则层”。

\begin{figure}[H]
\centering
\begin{tikzpicture}[
  node distance=1.15cm,
  box/.style={rectangle,draw,rounded corners,align=center,minimum width=2.6cm,minimum height=0.72cm},
  small/.style={rectangle,draw,rounded corners,align=center,minimum width=2.2cm,minimum height=0.62cm},
  arr/.style={-{Latex[length=2mm]},thick}
]
\node[box,fill=gray!10] (goal) {目标层\\综合评价目标 $G$};
\node[box,below left=1.3cm and 3.6cm of goal] (c1) {准则 $C_1$};
\node[box,below left=1.3cm and 1.1cm of goal] (c2) {准则 $C_2$};
\node[box,below right=1.3cm and 1.1cm of goal] (c3) {准则 $C_3$};
\node[box,below right=1.3cm and 3.6cm of goal] (c4) {准则 $C_4$};
\node[small,below=1.35cm of c1] (a1) {方案 $A_1$};
\node[small,below=1.35cm of c2] (a2) {方案 $A_2$};
\node[small,below=1.35cm of c3] (a3) {方案 $A_3$};
\node[small,below=1.35cm of c4] (a4) {方案 $A_m$};
\foreach \c in {c1,c2,c3,c4}{\draw[arr] (goal) -- (\c);}
\foreach \c in {c1,c2,c3,c4}{\foreach \a in {a1,a2,a3,a4}{\draw[arr,gray!55] (\c) -- (\a);}}
\node[draw,dashed,fit=(c1)(c4),inner sep=0.22cm,label=left:{准则层}]{};
\node[draw,dashed,fit=(a1)(a4),inner sep=0.22cm,label=left:{方案层}]{};
\end{tikzpicture}
\caption{AHP 三层评价结构图模板}
\label{fig:hierarchy}
\end{figure}

若仿照城市生态环境质量评价案例，可写成：目标层为“城市生态环境质量评价”；准则层包括自然环境、社会方面、经济方面、政策方面；二级指标可包括湿度、降水、PM2.5、废气、人均绿化、人口密度、人均 GDP、土地利用率、政策完善等。赵鑫等的案例中采用 4 个一级指标和 21 个二级指标，并将 AHP 与模糊综合评价结合使用。

\section{建模步骤与矩阵规范}
\subsection{Step1：明确评价问题}
先写清楚四件事：
\begin{enumerate}[leftmargin=2.2em]
  \item \textbf{评价目标：} 例如城市生态环境质量、供应商综合能力、项目风险等级。
  \item \textbf{评价对象或方案：} 记为 $A_1,A_2,\ldots,A_m$。
  \item \textbf{评价指标：} 记为 $C_1,C_2,\ldots,C_n$，必要时继续分解为 $C_{ij}$。
  \item \textbf{评分方向：} 明确哪些是效益型指标，哪些是成本型指标。
\end{enumerate}

\subsection{Step2：建立层次结构}
一般结构为
\begin{equation}
  G \longrightarrow \{C_1,C_2,\ldots,C_n\} \longrightarrow \{A_1,A_2,\ldots,A_m\}.
\end{equation}
若有子准则，则写成
\begin{equation}
  G \longrightarrow C_i \longrightarrow C_{ij} \longrightarrow A_k .
\end{equation}
建模论文中，层级结构图应放在方法部分或模型建立部分，用于说明评价体系来源和指标之间的从属关系。

\subsection{Step3：构造判断矩阵}
对准则层相对于目标层的重要性进行两两比较，得到判断矩阵
\begin{equation}
A=(a_{ij})_{n\times n}=\begin{pmatrix}
1 & a_{12} & \cdots & a_{1n}\\
1/a_{12} & 1 & \cdots & a_{2n}\\
\vdots & \vdots & \ddots & \vdots\\
1/a_{1n} & 1/a_{2n} & \cdots & 1
\end{pmatrix}.
\end{equation}
$a_{ij}$ 表示指标 $C_i$ 相对于 $C_j$ 的重要程度。常用 1--9 标度见表 \ref{tab:scale}。

\begin{table}[H]
\centering
\caption{AHP 1--9 标度表}
\label{tab:scale}
\begin{tabular}{cl}
\toprule
标度 & 含义\\
\midrule
1 & 两个因素同等重要\\
3 & 前者比后者稍微重要\\
5 & 前者比后者明显重要\\
7 & 前者比后者强烈重要\\
9 & 前者比后者极端重要\\
2,4,6,8 & 相邻判断之间的折中值\\
$1/3,1/5,\ldots,1/9$ & 若前者不如后者重要，则取倒数\\
\bottomrule
\end{tabular}
\end{table}

如果有多位专家，应先得到每位专家的判断矩阵 $A^{(r)}$，再用几何平均合成为群体判断矩阵：
\begin{equation}
  a_{ij}=\left(\prod_{r=1}^{q}a_{ij}^{(r)}\right)^{1/q} .
\end{equation}

\subsection{构造后先判断“能不能用”}
判断矩阵不能直接拿来算，至少要做以下检查：
\begin{enumerate}[leftmargin=2.2em]
  \item \textbf{矩阵合法性：} 是否为正互反矩阵，即 $a_{ij}>0$、$a_{ii}=1$、$a_{ij}=1/a_{ji}$。
  \item \textbf{指标口径一致性：} 同一层指标是否都服务于同一个上层目标，是否存在明显重复或遗漏。
  \item \textbf{评分方向一致性：} 方案评分前，效益型、成本型指标是否已经同向化。
  \item \textbf{一致性检验：} 判断矩阵的 $CR$ 是否小于阈值。常用经验标准为 $CR<0.10$。
\end{enumerate}

\subsection{Step4：依照指标对方案打分}
若方案层也采用两两比较，则对每个准则 $C_i$ 构造方案判断矩阵 $B_i$，得到方案在该准则下的局部权重：
\begin{equation}
  \bm v_i=(v_{i1},v_{i2},\ldots,v_{im})^T .
\end{equation}

若方案有客观数据，可先进行归一化。效益型指标越大越好：
\begin{equation}
  x'_{ik}=\frac{x_{ik}-\min_k x_{ik}}{\max_k x_{ik}-\min_k x_{ik}} .
\end{equation}
成本型指标越小越好：
\begin{equation}
  x'_{ik}=\frac{\max_k x_{ik}-x_{ik}}{\max_k x_{ik}-\min_k x_{ik}} .
\end{equation}
若 $\max x_{ik}=\min x_{ik}$，说明该指标不能区分方案，可赋为同一得分或剔除该指标。

\subsection{Step5：计算权重、填表并求最终得分}
常用的特征向量法如下：
\begin{enumerate}[leftmargin=2.2em]
  \item 求判断矩阵 $A$ 的最大特征根 $\lambda_{\max}$；
  \item 取其对应的正特征向量 $\bm w$；
  \item 归一化，使 $\sum_i w_i=1$。
\end{enumerate}

若采用客观评分矩阵 $X'=(x'_{ik})_{n\times m}$，则方案 $A_k$ 的总分为
\begin{equation}
  S_k=\sum_{i=1}^{n} w_i x'_{ik} .
\end{equation}
若方案层也用 AHP 两两比较，则将各准则下的方案局部权重组成矩阵
\begin{equation}
V=\begin{pmatrix}
 v_{11}&v_{12}&\cdots&v_{1m}\\
 v_{21}&v_{22}&\cdots&v_{2m}\\
 \vdots&\vdots&\ddots&\vdots\\
 v_{n1}&v_{n2}&\cdots&v_{nm}
\end{pmatrix},
\end{equation}
综合得分向量为
\begin{equation}
  \bm S=\bm w^T V .
\end{equation}

论文中的结果表建议按表 \ref{tab:result-template} 写。
\begin{table}[H]
\centering
\caption{AHP 权重与得分结果表模板}
\label{tab:result-template}
\begin{tabular}{cccccc}
\toprule
指标 & 局部权重 & 方案 $A_1$ & 方案 $A_2$ & $\cdots$ & 方案 $A_m$\\
\midrule
$C_1$ & $w_1$ & $x'_{11}$ & $x'_{12}$ & $\cdots$ & $x'_{1m}$\\
$C_2$ & $w_2$ & $x'_{21}$ & $x'_{22}$ & $\cdots$ & $x'_{2m}$\\
$\vdots$ & $\vdots$ & $\vdots$ & $\vdots$ & $\ddots$ & $\vdots$\\
$C_n$ & $w_n$ & $x'_{n1}$ & $x'_{n2}$ & $\cdots$ & $x'_{nm}$\\
\midrule
综合得分 & -- & $S_1$ & $S_2$ & $\cdots$ & $S_m$\\
排序 & -- &  &  &  & \\
\bottomrule
\end{tabular}
\end{table}

\subsection{Step6：一致性检验与层次总排序一致性}
判断矩阵一致性指标为
\begin{equation}
  CI=\frac{\lambda_{\max}-n}{n-1} .
\end{equation}
随机一致性指标 $RI$ 由矩阵阶数决定，常用取值见表 \ref{tab:ri}。
\begin{table}[H]
\centering
\caption{随机一致性指标 RI 常用表}
\label{tab:ri}
\begin{tabular}{c|cccccccccc}
\toprule
$n$ & 1 & 2 & 3 & 4 & 5 & 6 & 7 & 8 & 9 & 10\\
\midrule
$RI$ & 0 & 0 & 0.58 & 0.90 & 1.12 & 1.24 & 1.32 & 1.41 & 1.45 & 1.49\\
\bottomrule
\end{tabular}
\end{table}
一致性比率为
\begin{equation}
  CR=\frac{CI}{RI} .
\end{equation}
当 $n=1$ 或 $n=2$ 时，可认为矩阵总是一致的；当 $n\ge 3$ 时，一般要求 $CR<0.10$。若 $CR\ge 0.10$，说明专家判断矛盾较大，应重新比较相近指标或重新组织专家打分。

对于多层结构，除了每个局部判断矩阵要通过一致性检验，还要做层次总排序一致性检验。设准则层权重为 $w_i$，各准则下方案判断矩阵的局部一致性指标为 $CI_i$、随机一致性指标为 $RI_i$，则
\begin{equation}
  CI_{total}=\sum_{i=1}^{n}w_iCI_i,
  \qquad
  RI_{total}=\sum_{i=1}^{n}w_iRI_i,
\end{equation}
\begin{equation}
  CR_{total}=\frac{CI_{total}}{RI_{total}} .
\end{equation}
若 $CR_{total}<0.10$，则层次总排序通过一致性检验，最终排序可以采用；否则需要调整判断矩阵。

\section{一个小型算例：三指标、三方案}
设评价目标为“选择最优方案”，准则为成本 $C_1$、质量 $C_2$、风险 $C_3$。专家给出准则判断矩阵
\begin{equation}
A=\begin{pmatrix}
1 & 3 & 5\\
1/3 & 1 & 2\\
1/5 & 1/2 & 1
\end{pmatrix}.
\end{equation}
计算可得准则权重近似为
\begin{equation}
  \bm w=(0.6483,\;0.2297,\;0.1220)^T,
\end{equation}
$\lambda_{\max}\approx 3.0037$，$CI\approx0.0018$，$CR\approx0.0032<0.10$，通过一致性检验。

若三个方案在三个指标上的标准化得分矩阵为
\begin{equation}
X'=\begin{pmatrix}
0.90&0.70&0.40\\
0.60&0.85&0.75\\
0.80&0.50&0.90
\end{pmatrix},
\end{equation}
则
\begin{equation}
\bm S=\bm w^T X'=(0.8189,\;0.7100,\;0.5414).
\end{equation}
因此排序为 $A_1>A_2>A_3$。

\section{结论评价的论文写法}
最终结论建议按“排序结果--原因解释--可靠性说明--对策建议”四段写。

\subsection{通用结论模板}
在既定指标体系和专家判断条件下，方案 $A_i$ 的综合得分最高，说明其在主要评价指标上整体表现最优。进一步观察权重可知，影响总排序的关键指标为 $C_p$、$C_q$ 等，说明本问题中这些因素对评价结果具有主导作用。由于所有判断矩阵均通过一致性检验，且层次总排序一致性比率小于 0.10，故所得排序具有较好的逻辑一致性和可解释性。

若某指标权重明显偏高，应结合实际解释其合理性。例如生态环境评价中，政策完善、废气治理、人均绿化等指标权重较高，可以解释为城市生态环境不仅受自然条件影响，也受治理政策、经济结构和社会发展水平影响。

\subsection{结合城市生态环境案例的评价写法}
以城市生态环境质量评价为例，赵鑫等的研究建立了自然环境、社会方面、经济方面和政策方面四个一级指标，以及 21 个二级指标。其 AHP 计算结果显示，一级指标中社会方面权重最大，为 0.342；自然环境权重为 0.316；复合权重中政策完善权重最大，为 0.1007，废气次之，为 0.0752。进一步结合模糊综合评价后，该城市生态环境质量评价值约为 0.5978，属于中等水平。这个结论的写法重点不是只报数值，而是说明：当前城市整体生态环境还有提升空间，治理时应优先关注权重较高、且评分不高的指标。

\subsection{局限性说明}
AHP 的局限性也应在论文中交代：
\begin{enumerate}[leftmargin=2.2em]
  \item 结果依赖指标体系，若指标遗漏或重复，权重会失真。
  \item 专家判断具有主观性，应说明专家来源、人数和专业背景。
  \item 1--9 标度只能近似表达偏好强度，不宜过度解释为精确物理量。
  \item 加权和模型默认指标之间可补偿，若某些指标具有“一票否决”性质，需要另设约束或阈值。
  \item 最好补充敏感性分析，检验权重小幅变化时排序是否稳定。
\end{enumerate}

\section{Python 代码模板}
下面代码可以直接用于 AHP 建模。它包含权重计算、一致性检验、群体判断矩阵合成、客观指标归一化和最终得分排序。

\begin{lstlisting}
import numpy as np

RI_TABLE = {
    1: 0.00, 2: 0.00, 3: 0.58, 4: 0.90, 5: 1.12,
    6: 1.24, 7: 1.32, 8: 1.41, 9: 1.45, 10: 1.49,
    11: 1.51, 12: 1.53, 13: 1.56, 14: 1.57, 15: 1.59
}


def check_reciprocal(A, tol=1e-8):
    """Check whether A is a positive reciprocal matrix."""
    A = np.asarray(A, dtype=float)
    if A.ndim != 2 or A.shape[0] != A.shape[1]:
        raise ValueError("A must be a square matrix.")
    if np.any(A <= 0):
        raise ValueError("All entries of A must be positive.")
    n = A.shape[0]
    if not np.allclose(np.diag(A), np.ones(n), atol=tol):
        raise ValueError("Diagonal entries must be 1.")
    if not np.allclose(A * A.T, np.ones_like(A), atol=tol):
        raise ValueError("A must satisfy a_ij = 1/a_ji.")
    return True


def ahp_weights(A):
    """Return weights, lambda_max, CI, CR for a judgment matrix."""
    A = np.asarray(A, dtype=float)
    check_reciprocal(A)
    n = A.shape[0]

    eigvals, eigvecs = np.linalg.eig(A)
    idx = np.argmax(eigvals.real)
    lambda_max = eigvals[idx].real
    w = eigvecs[:, idx].real

    # The principal eigenvector can be multiplied by -1.
    w = np.abs(w)
    w = w / w.sum()

    if n <= 2:
        CI, CR = 0.0, 0.0
    else:
        CI = (lambda_max - n) / (n - 1)
        RI = RI_TABLE.get(n)
        if RI is None:
            raise ValueError("RI value is not available for this n.")
        CR = CI / RI if RI != 0 else 0.0
    return w, lambda_max, CI, CR


def group_geometric_mean(matrices):
    """Aggregate several experts' judgment matrices by geometric mean."""
    mats = np.asarray(matrices, dtype=float)
    if mats.ndim != 3:
        raise ValueError("Input shape should be (q, n, n).")
    for A in mats:
        check_reciprocal(A)
    return np.prod(mats, axis=0) ** (1.0 / mats.shape[0])


def normalize_scores(X, types):
    """
    Normalize objective data.
    types[i] = 'benefit' means larger is better.
    types[i] = 'cost' means smaller is better.
    Rows are criteria, columns are alternatives.
    """
    X = np.asarray(X, dtype=float)
    Xn = np.zeros_like(X)
    for i, t in enumerate(types):
        row = X[i]
        xmin, xmax = row.min(), row.max()
        if np.isclose(xmax, xmin):
            Xn[i] = 1.0
        elif t == "benefit":
            Xn[i] = (row - xmin) / (xmax - xmin)
        elif t == "cost":
            Xn[i] = (xmax - row) / (xmax - xmin)
        else:
            raise ValueError("types must be 'benefit' or 'cost'.")
    return Xn


def rank_by_ahp(criteria_matrix, score_matrix):
    """Use criteria weights and standardized score matrix to rank alternatives."""
    w, lambda_max, CI, CR = ahp_weights(criteria_matrix)
    scores = w @ np.asarray(score_matrix, dtype=float)
    order = np.argsort(-scores)  # descending order
    return {
        "weights": w,
        "lambda_max": lambda_max,
        "CI": CI,
        "CR": CR,
        "scores": scores,
        "ranking_index": order
    }


if __name__ == "__main__":
    A = np.array([
        [1,   3,   5],
        [1/3, 1,   2],
        [1/5, 1/2, 1]
    ], dtype=float)

    X_std = np.array([
        [0.90, 0.70, 0.40],
        [0.60, 0.85, 0.75],
        [0.80, 0.50, 0.90]
    ], dtype=float)

    result = rank_by_ahp(A, X_std)
    print("weights =", np.round(result["weights"], 4))
    print("lambda_max =", round(result["lambda_max"], 4))
    print("CI =", round(result["CI"], 4))
    print("CR =", round(result["CR"], 4))
    print("scores =", np.round(result["scores"], 4))
    print("ranking =", result["ranking_index"] + 1)
\end{lstlisting}

\section{写进论文时的规范性检查清单}
\begin{enumerate}[leftmargin=2.2em]
  \item 是否解释了为什么这是评价类问题，为什么适合 AHP？
  \item 是否给出了目标层、准则层、方案层结构图？
  \item 判断矩阵是否满足 $a_{ij}=1/a_{ji}$、$a_{ii}=1$？
  \item 是否说明 1--9 标度含义和专家打分来源？
  \item 是否计算了 $\lambda_{\max}$、$CI$、$RI$、$CR$？
  \item 若有多层结构，是否做了层次总排序一致性检验？
  \item 方案评分是否已经无量纲化、同向化？
  \item 结论是否既有排名，又解释了关键影响指标？
  \item 是否说明局限性和敏感性分析？
\end{enumerate}

\begin{thebibliography}{9}
\bibitem{Saaty1987} Saaty, R. W. The analytic hierarchy process—what it is and how it is used. \emph{Mathematical Modelling}, 1987, 9(3--5): 161--176. DOI: 10.1016/0270-0255(87)90473-8.
\bibitem{Salomon2024} Salomon, V. A. P. Consistency Improvement in the Analytic Hierarchy Process. \emph{Mathematics}, 2024, 12(6): 828.
\bibitem{RI} ScienceDirect Topics. Consistency Index: Estimation of consistency ratio and random index table.
\bibitem{Steps} 1000minds. What is the Analytic Hierarchy Process (AHP)? Structure hierarchy, pairwise comparison, weights, alternative evaluation, weighted-sum ranking.
\bibitem{Zhao2022} 赵鑫，孙春花，沈贤. 基于层次分析法的城市生态环境质量评价[J]. 中国资源综合利用, 2022, 40(5): 163--166.
\end{thebibliography}

\end{document}
